\documentclass[10pt,a4paper]{article}

\usepackage[margin=1.75cm]{geometry}
\usepackage[T1]{fontenc}
\usepackage[utf8]{inputenc}
\usepackage{lmodern}
\usepackage{amsmath,amssymb}
\usepackage{xcolor}
\usepackage{microtype}
\usepackage{listings}
\usepackage[hidelinks]{hyperref}

\definecolor{codegray}{HTML}{F5F6F8}
\definecolor{rulegray}{HTML}{D6DAE0}

\lstdefinestyle{pythoncompact}{
  language=Python,
  basicstyle=\ttfamily\scriptsize,
  backgroundcolor=\color{codegray},
  frame=single,
  rulecolor=\color{rulegray},
  showstringspaces=false,
  columns=fullflexible,
  keepspaces=true,
  aboveskip=0.4em,
  belowskip=0.4em
}

\newcommand{\E}{\mathbb{E}}
\newcommand{\1}{\mathbf{1}}
\newcommand{\argmax}{\operatorname*{arg\,max}}

\setlength{\parindent}{0pt}
\setlength{\parskip}{3pt}
\setlength{\abovedisplayskip}{5pt}
\setlength{\belowdisplayskip}{5pt}

\title{\vspace{-1.4em}\textbf{Requirement 4: Why the Observed Regret is Linear and How to Correct the Baseline}}
\author{OLA Project -- Technical Note}
\date{}

\begin{document}
\maketitle
\vspace{-1.6em}

\begin{abstract}
\noindent Requirement 4 correctly implements a non-stationary multiple-campaign auction with a shared budget: the budget plot is coherent, and cumulative cost approaches the total budget $B$ near the end of the horizon. The apparent issue is that cumulative regret can look almost linear. This does not necessarily indicate a failure of Sliding-Window Combinatorial-UCB, CUSUM Combinatorial-UCB, or the Primal-Dual baseline. The main cause is the benchmark currently used in the code: a dynamic realized clairvoyant, or prophet oracle, that knows the realized competing bids round by round. This adds a persistent information advantage. A more meaningful diagnostic baseline is a piecewise expected clairvoyant computed from the stationary distribution inside each interval.
\end{abstract}

\section{Theoretical Setting}
The horizon has $T$ rounds, $N$ campaigns, one shared budget $B$, and average budget
\[
  \rho = B/T .
\]
Requirement 4 is piecewise stationary: the horizon is divided into intervals $I_1,\ldots,I_S$. Inside interval $I_s$, the distribution $P_s$ of the highest competing bids is fixed; across intervals, it changes. Let $m_{i,t}$ be the highest competing bid for campaign $i$ at time $t$, and let $v_i$ be its value. For campaign $i$ and bid $b$,
\[
  u_t(i,b)=(v_i-b)\1\{b\ge m_{i,t}\},\qquad
  c_t(i,b)=b\1\{b\ge m_{i,t}\}.
\]
For a joint action $a=(b_1,\ldots,b_N)$,
\[
  U_t(a)=\sum_i u_t(i,b_i),\qquad
  C_t(a)=\sum_i c_t(i,b_i).
\]
The realized budget constraint is
\[
  \sum_{t=1}^T C_t(a_t)\le B,
\]
while the stationary LP view uses the per-round expected constraint
\[
  \E[C_t(a_t)]\le \rho .
\]
The algorithms in Req4 estimate these quantities online and optimize over feasible joint actions subject to shared-budget and conflict-graph constraints.

\section{Current Regret Implementation}
The current regret uses a dynamic realized clairvoyant. In simplified form:
\begin{lstlisting}[style=pythoncompact]
dynamic_oracle = compute_dynamic_clairvoyant_realized(env, budget)
regret = np.cumsum(dynamic_oracle.utilities - learner.utilities)
\end{lstlisting}
In the local code, this corresponds to a dynamic LP over the entire realized sequence \texttt{m\_seq}. It knows every $m_{i,t}$ before selecting the oracle action, then compares the learner against the oracle's realized utility. Mathematically,
\[
  R_T^{\mathrm{prophet}}
  =
  \sum_{t=1}^T U_t(a_t^{\mathrm{prophet}})
  -
  \sum_{t=1}^T U_t(a_t),
\]
where
\[
  a_t^{\mathrm{prophet}}
  =
  \argmax_{a\in\mathcal A_t(B)} U_t(a)
\]
is chosen after observing the realized vector $(m_{1,t},\ldots,m_{N,t})$. This is stronger than knowing the interval distribution $P_s$: it is an ex-post selector over the random auction outcomes.

\section{Why This Creates a Linear Component}
Even within a stationary interval, $m_{i,t}$ is freshly drawn each round. A learner, and even a fair distribution-aware clairvoyant, may know $P_s$ but cannot know the realized $m_t$ before bidding. The prophet can condition on $m_t$. Therefore,
\[
  \E\!\left[\max_{a\in\mathcal A} U(a;m)\right]
  \ge
  \max_{\pi\in\Pi}\E\!\left[U(a\sim\pi;m)\right],
\]
with strict inequality whenever different realizations make different actions optimal. Define the interval information premium
\[
\Delta_s=
\E_{m\sim P_s}[U(a^{\mathrm{prophet}}(m);m)]
-
\max_{\pi\in\Pi}\E_{m\sim P_s}[U(a\sim\pi;m)] .
\]
If $\Delta_s>0$, a learner that has already identified the best feasible expected policy still loses about $\Delta_s$ utility per round against the prophet. Hence
\[
  R_T^{\mathrm{prophet}}
  =
  R_T^{\mathrm{learn}}
  +
  \sum_{s=1}^S |I_s|\Delta_s
  \quad
  \text{(up to budget-coupling effects).}
\]
The second term is linear in $T$. Thus a near-linear regret curve can be a baseline artifact, not proof that the algorithms have linear learning regret.

\section{Correct Baseline}
The recommended benchmark is a piecewise expected clairvoyant. For each interval $I_s$, compute
\[
  \bar U_s(a)=\E_{m\sim P_s}[U(a;m)],\qquad
  \bar C_s(a)=\E_{m\sim P_s}[C(a;m)].
\]
Then solve the shared-budget LP
\[
\begin{aligned}
  \max_{x_{s,a}\ge0}\quad
  &\sum_s |I_s|\sum_a x_{s,a}\bar U_s(a)\\
  \text{s.t.}\quad
  &\sum_s |I_s|\sum_a x_{s,a}\bar C_s(a)\le B,\\
  &\sum_a x_{s,a}=1,\qquad s=1,\ldots,S .
\end{aligned}
\]
The resulting regret is
\[
  R_T^{\mathrm{pw-exp}}
  =
  \sum_s |I_s|V_s^\star
  -
  \sum_{t=1}^T U_t(a_t).
\]
This oracle may know the changepoints and the distribution of each stationary segment, but not the realized competing bids before bidding. It is therefore a stochastic non-stationary clairvoyant rather than a prophet.

\section{Implementation Sketch}
The learners do not need to change first. Add a second benchmark and report the realized dynamic oracle only as a prophet upper bound:
\begin{lstlisting}[style=pythoncompact]
pw_oracle = compute_piecewise_expected_clairvoyant(
    intervals=env.intervals,
    bid_sets=env.bid_sets,
    values=env.values,
    budget=env.budget,
    conflict_edges=env.conflict_edges,
)
regret = np.cumsum(pw_oracle.expected_per_round - learner.utilities)
\end{lstlisting}
In the synthetic environment, $\bar U_s$ and $\bar C_s$ can be computed from the known generator parameters, or by large independent Monte Carlo samples within each interval.

\section{Conclusion}
The budget plot suggests that Req4's shared-budget mechanism is working: cumulative cost approaches $B$ near $T$ rather than exhausting the budget too early or leaving it largely unused. The near-linear regret is consistent with comparing online algorithms to a prophet that observes realized bids round by round. Replacing the main benchmark with a piecewise expected clairvoyant separates genuine adaptation regret from the artificial linear information premium introduced by the dynamic realized oracle.

\end{document}
